\section{A budget for pivotal labels}
\label{sec:pivotality}

We associate to each conditioned function a quantity that controls both its
influences and its acceptance probabilities at every bias.  The single-label
tests bound the expectation of this quantity under the common distribution
from Section~\ref{sec:prefix}.
A single test vertex may have several cuts that force its deletion, but its
cost is paid only once.  Accordingly, we record whether a label is pivotal
somewhere along a permutation, rather than the number of witnessing prefixes.

\begin{definition}[Pivotality budget]
\label{def:pivotal-budget}
Let $S$ be a nonempty finite set, and let
$f:2^S\to\{0,1\}$ be monotone with $f(\varnothing)=0$ and $f(S)=1$.
For each $j\in S$, let $\sigma_j$ be a uniform permutation of
$S\setminus\{j\}$.  Write $R_a(\sigma_j)$ for the set of its first $a$
elements, where $0\le a\le |S|-1$.  Define
\[
 B(f)=\sum_{j\in S}
 \Pr_{\sigma_j}\!\left[
 \begin{gathered}
   \text{there exists }a\in\{0,\ldots,|S|-1\}\text{ such that}\\
   f(R_a(\sigma_j))=0,\quad
   f(R_a(\sigma_j)\cup\{j\})=1
 \end{gathered}
 \right].
\]
Thus both the empty prefix and the full prefix $S\setminus\{j\}$ are
allowed.  A permutation contributes only once for a given $j$, even when
several of its prefixes witness pivotality.  In particular,
$0\le B(f)\le |S|$.
\end{definition}

\begin{lemma}[Arrival-time inequalities]
\label{lem:pivotal-inequalities}
For every function in Definition~\ref{def:pivotal-budget} and every
$0<p<1$,
\[
 I_p(f)\le B(f),
 \qquad
 \mathbb E_{R\sim\mu_p^S} f(R)\le pB(f).
\]
\end{lemma}

\begin{proof}
Give each $j\in S$ an independent uniform arrival time $T_j\in[0,1]$.
With probability one all these times are distinct.  For each $j$, let
$\sigma_j$ be the order in which the other elements arrive, and let
$\mathcal A_j$ be the event that some prefix of $\sigma_j$ witnesses
$j$'s pivotality in Definition~\ref{def:pivotal-budget}.  The permutation
$\sigma_j$ is uniform, so
\[
 \sum_{j\in S}\Pr(\mathcal A_j)=B(f).
\]
Moreover, $\mathcal A_j$ is determined entirely by the arrival times of
elements other than $j$, and hence is independent of $T_j$.

For the influence inequality, put
\[
 R_{-j}(p)=\{k\in S\setminus\{j\}:T_k\le p\}.
\]
This set has law $\mu_p^{S\setminus\{j\}}$ and is a prefix of
$\sigma_j$.  By monotonicity, the event defining $I_{p,j}(f)$ is exactly
\[
 \{f(R_{-j}(p))=0,\ f(R_{-j}(p)\cup\{j\})=1\},
\]
which is contained in $\mathcal A_j$.  Summing
$I_{p,j}(f)\le\Pr(\mathcal A_j)$ over $j$ proves $I_p(f)\le B(f)$.

For the acceptance inequality, put $R(p)=\{j\in S:T_j\le p\}$, which
has law $\mu_p^S$.  Along the complete arrival sequence, the value of
$f$ starts at zero, ends at one, and never decreases.  There is therefore
a unique arriving element at which it changes from zero to one.  If
$f(R(p))=1$ and this element is $j$, then $T_j\le p$, and the elements
arriving before $j$ form a prefix of $\sigma_j$ witnessing
$\mathcal A_j$.  Consequently, up to the null event of tied arrival times,
\[
 \{f(R(p))=1\}
 \subseteq\bigcup_{j\in S}\bigl(\{T_j\le p\}\cap\mathcal A_j\bigr).
\]
A union bound and the independence established above give
\[
 \mathbb E_{R\sim\mu_p^S} f(R)
 \le\sum_{j\in S}\Pr(T_j\le p,\mathcal A_j)
 =p\sum_{j\in S}\Pr(\mathcal A_j)
 =pB(f).
\]
Only the possible-witness event $\mathcal A_j$ is asserted to be
independent of $T_j$; no independence of the actual switching element is
needed.
\end{proof}

Return to the fixed feedback vertex set $F$ of cost at most $D$ and its
fixed topological order.  For the functions in
Definition~\ref{def:conditioned-functions}, set
\[
 B_w=B(f_w)\qquad(w\in U\cup V).
\]
These variables depend on the shared data $\mathcal D$ and the named
vertex $w$.  Their definition does not involve a fresh tested edge or a
fresh bias index.  Lemma~\ref{lem:pivotal-inequalities} gives, for every
realization of these data and every $0<p<1$,
\[
 I_p(f_w)\le B_w,
 \qquad h_w(p)\le pB_w.
\]

\begin{proposition}[Single-label test budget]
\label{prop:pivotal-budget}
Under the common law of $\mathcal D$ and an independent edge
$e=(u,v)\sim\nu$,
\begin{equation}
 \mathbb E B_u\le D+1,
 \qquad
 \mathbb E B_v\le D+1.
 \label{eq:pivotal-budget}
\end{equation}
\end{proposition}

\begin{proof}
Fix either side of the game and sample $w$ independently of
$\mathcal D$ from that side's endpoint marginal.  Its alphabet size
$m_w$ is constant on the chosen side.  This is exactly the marginal law
of the shared data and endpoint in the corresponding single-label test
family of Definition~\ref{def:tests}.

First fix $\mathcal D,w$ for which the extension $\mathbf s w$ is good.
The conditioned function $f_w$ then uses the transition cell $C_{r_*}$.
The remaining randomness of a single-label test consists of a uniform
label $j_0\in\Lambda_w$ and an independent uniform permutation $\sigma$
of $\Lambda_w$.  Its star set is
\[
 Z=\Lambda_{\mathbf s}\times\{j_0\},
\]
and its ordered cells are $C_\ell\times\{j\}$, first ordered by
$\ell$ and then by $\sigma$.  Suppose some prefix $R$ of $\sigma$ with
$j_0$ omitted satisfies
\[
 f_w(R)=0,
 \qquad f_w(R\cup\{j_0\})=1.
\]

Choose a cut among the refined cells belonging to $C_{r_*}$, including
either boundary of this block if necessary, so that precisely the
nonstar labels in $R$ have been passed within that block.  Such a cut
exists because deleting $j_0$ from the sequence $\sigma$ preserves the
order of every other label.  For $R=\varnothing$ one may cut before the
block, and for $R=\Lambda_w\setminus\{j_0\}$ one may cut after the
block.  Let $E$ be the union of all refined cells before the cut, with
all stars removed.  Then exactly
\[
 E=\bigl((C_{<r_*}\times\Lambda_w)\setminus Z\bigr)
       \cup(C_{r_*}\times R).
\]
In particular,
\[
 E\subseteq E_w^{r_*}(R),
 \qquad
 E\cup Z\supseteq E_w^{r_*}(R\cup\{j_0\}).
\]
The second inclusion restores the stars in all earlier prefix cells;
the additional stars in later cells only enlarge the set.  Monotonicity
of $g_{\mathbf s w}$ and the definition of $f_w$ now imply
\[
 g_{\mathbf s w}(E)=0,
 \qquad g_{\mathbf s w}(E\cup Z)=1.
\]
Lemma~\ref{lem:wildcard-forcing} therefore forces this test vertex to
belong to $F$.

This argument uses the ordered partition with its empty cells retained;
the set identities remain valid when some of those cells are empty.
They also hold on the formal product label domain when $w$ or another
game vertex occurs repeatedly in the tuple.  No independence of labels
in repeated slots is assumed, and the use of $f_w$ rests only on the
prefix identifications of Lemma~\ref{lem:compatibility}.

For each fixed $j_0$, deleting $j_0$ from a uniform permutation of
$\Lambda_w$ gives a uniform permutation of
$\Lambda_w\setminus\{j_0\}$.  Averaging the forced-deletion event over
the uniform choice of $j_0$ therefore gives, whenever $\mathbf s w$ is
good,
\[
 \Pr[\text{sampled test vertex}\in F\mid\mathcal D,w]
 \ge \frac{B_w}{m_w}.
\]
Here a test is charged once if a witnessing prefix exists, regardless
of how many prefixes witness the same label's pivotality.

Let $\mathcal G_w$ denote the event that $\mathbf s w$ is good.
The total mass of the chosen test family is $m_w$, and its vertices
carry their sampling probabilities times this mass.  Thus, on taking
expectations, the cost spent by $F$ in this family is at least
$\mathbb E[B_w\mathbf 1_{\mathcal G_w}]$.  Since the total cost of
$F$ is at most $D$,
\[
 \mathbb E[B_w\mathbf 1_{\mathcal G_w}]\le D.
\]
On the complementary event the fallback function is still monotone,
with bottom zero and top one, so Definition~\ref{def:pivotal-budget}
gives $B_w\le m_w$.  The sampled-extension bound from
Equation~\eqref{eq:sampled-extension-bad} gives
$\Pr(\mathcal G_w^c)\le\rho$; equivalently, one may apply
Lemma~\ref{lem:goodness} after fixing the sampled vertex tuple, which
is independent of $\psi$.  As $m_w\le m$ and $\rho\le1/m$,
\[
 \mathbb E[B_w\mathbf 1_{\mathcal G_w^c}]
 \le m_w\rho\le1.
\]
Adding the two bounds proves $\mathbb E B_w\le D+1$ for either
endpoint marginal, and hence proves Equation~\eqref{eq:pivotal-budget}.
\end{proof}
